\section{Related Work}

\subsection{Cross-Modal Grounding and Temporal RVOS}
A central line of RVOS research couples language grounding with temporal mask prediction. OnlineRefer~\cite{wu2023onlinerefer} propagates object queries in an online setting, while DsHmp~\cite{he2024dshmp} explicitly separates static appearance from hierarchical motion cues to better model motion-dependent expressions. More recently, ReferDINO~\cite{liang2025referdino} transfers pretrained visual-grounding representations to RVOS through temporal enhancement, grounding-guided mask decoding, and query pruning. These methods differ substantially in architecture, but they share a common formulation: the representation used to identify the referent is tightly coupled to the representation used for dense mask prediction and temporal propagation. This coupling is effective when grounding and tracking succeed together, yet it makes a missed target hypothesis difficult to distinguish from an incorrect downstream selection. CRF-SAM3 takes a different view by making the object hypothesis space explicit: language and video first determine what should be recalled, while final referent selection is postponed until candidate tubes have been constructed.

\subsection{MLLMs and Foundation Video Segmenters}
MLLM-based methods move more of the semantic burden from task-specific cross-modal encoders to pretrained language--vision models. VideoLISA~\cite{bai2024videolisa} extends segmentation-token-based reasoning to videos, and Sa2VA~\cite{yuan2025sa2va} connects an MLLM with SAM2 so that language-model outputs can guide dense image and video segmentation. VRS-HQ~\cite{gong2025vrshq} further strengthens temporal information encoded by segmentation tokens. A related direction adapts strong video segmenters themselves: SAMWISE~\cite{cuttano2024samwise} injects language and explicit temporal reasoning into SAM2 while preserving its pretrained segmentation and memory capabilities. GLUS~\cite{lin2025glus} instead addresses the tension between global referring semantics and local video tracking by jointly using sparse context frames, continuous query frames, and a VOS memory bank. StreamingRVOS~\cite{zhang2026streamingrvos} targets online inference by recycling semantic embeddings over time and re-invoking multimodal perception when mask consistency degrades.

These works mainly ask how language semantics should be fused into, propagated through, or repeatedly supplied to a video segmentation system. SAM3 provides another useful interface through promptable concept segmentation~\cite{carion2025sam3}: a compact open-vocabulary concept can be used to discover and track matching instances. We use this capability in a deliberately restricted way. Rather than treating the foundation segmenter as the final resolver of a compositional referring expression, CRF-SAM3 uses a visually grounded concept to \emph{recall} plausible object tubes and then restores the complete expression for candidate-level verification. Thus, the compact concept serves object discovery, whereas fine-grained action, relation, and temporal semantics remain available for later discrimination.

\subsection{Decoupled Identification and Trajectory-Level Reasoning}
A growing body of work questions whether target identification and temporal segmentation must be solved by one tightly entangled representation. FindTrack~\cite{cho2025findtrack} first identifies the target on an informative frame and then delegates temporal propagation to a dedicated tracker, separating ``find'' from ``track.'' DeRVOS~\cite{cheng2026dervos} pushes this idea to the trajectory level: temporally consistent instance trajectories are generated with a pretrained trajectory model and subsequently aligned with multimodal representations for trajectory selection. Refer-Agent~\cite{jiang2026referagent} explores a different form of decomposition without task-specific fine-tuning, using collaborative agents, adaptive frame selection, focused visual layouts, and iterative reasoning--reflection to refine the referred target.

CRF-SAM3 is closest to this family in spirit, but the boundary between its stages is different. We do not first commit to a single target and then propagate it, nor do we reduce RVOS only to matching language against a fixed set of pretrained trajectories. Instead, we learn both interfaces around an explicit SAM3 candidate space: an expression-conditioned grounder determines a recall concept before candidate generation, and a task-supervised verifier learns competition among the recalled tubes afterward. This explicit interface is also diagnostic: a target can fail because it was never recalled, or because a valid candidate was recalled but ranked incorrectly. The experiments analyze these two failure modes separately rather than attributing both to the final mask predictor.
